% Transcription of Nguyen Van Khue and Nguyen To Nhu (1981), pp. 285--291.
% Source provenance and original misprints: TRANSCRIPTION-NOTES.md.
\documentclass[11pt,a4paper,leqno,twoside]{article}
\usepackage[T2A,T1]{fontenc}
\usepackage[utf8]{inputenc}
\usepackage[russian,english]{babel}
\usepackage{amsmath,amssymb}
\usepackage[margin=22mm,headheight=15pt]{geometry}
\usepackage{fancyhdr,enumitem}
\pagestyle{fancy}
\fancyhf{}
\fancyhead[LE,RO]{\thepage}
\fancyhead[CE]{Nguyen Van Khue, Nguyen To Nhu}
\fancyhead[CO]{\textit{Two Extensors of Metrics}}
\renewcommand{\headrulewidth}{0.3pt}
\setlength{\parindent}{1.5em}
\setlength{\parskip}{2pt}
\setlength{\emergencystretch}{1em}
\newcommand{\Metr}{\operatorname{Metr}}
\newcommand{\Pm}{\operatorname{Pm}}
\newcommand{\Lip}{\operatorname{Lip}}
\newcommand{\conv}{\operatorname{conv}}
\newcommand{\norm}[1]{\lVert#1\rVert}
\newcommand{\restr}{\vert}
\newcommand{\proofpart}[1]{\par\medskip\noindent Proof of (#1).\enspace}
\newenvironment{conditions}[1]{\begin{enumerate}[label=(#1-\roman*),leftmargin=4.3em,itemsep=1pt,topsep=3pt,parsep=0pt]}{\end{enumerate}}
\begin{document}
\setcounter{page}{285}
\thispagestyle{plain}
\noindent\begin{tabular}{@{}c@{}}
\textsc{Bulletin de l'Acad\'emie}\\
\textsc{Polonaise des Sciences}\\
S\'erie des sciences math\'ematiques\\
Vol. XXIX, No. 5--6, 1981
\end{tabular}
\vspace{7mm}
\begin{flushright}\textit{GENERAL TOPOLOGY}\end{flushright}
\begin{center}
{\Large Two Extensors of Metrics}\par\medskip
by\par\medskip
NGUYEN VAN KHUE and NGUYEN TO NHU\par\medskip
\textit{Presented by K. BORSUK on November 13, 1979}
\end{center}
\medskip
\begingroup\small
\noindent\textbf{Summary.} By $\Metr(Z)$ (resp.\ $\Metr_c(Z)$) we denote the family of all compatible metrics (resp.\ complete compatible metrics) on a given metrizable space $Z$. Let $A$ be a closed subset of a metrizable space $X$. In this note we construct two extensors of metrics $\Phi_1,\Phi_2:\Metr(A)\to\Metr(X)$ satisfying the conditions
\begin{enumerate}[label=(\roman*),leftmargin=3.5em,itemsep=0pt,topsep=2pt]
\item $\Phi_i(d)\restr A=d$ for $d\in\Metr(A)$ and $i=1,2$.
\item If $X$ is complete-metrizable, then
\[\Phi_i(\Metr_c(A))\subset\Metr_c(X)\quad\text{for }i=1,2.\]
\item $\Phi_1$ is a Lipschitz map.
\item If $d_1\leq d_2$ then $\Phi_2(d_1)\leq\Phi_2(d_2)$ for $d_1,d_2\in\Metr(A)$.
\item If $d_n(x,y)\to d_0(x,y)$ for $x,y\in A$, then
\[\Phi_2(d_n)(x,y)\to\Phi_2(d_0)(x,y)\quad\text{for }x,y\in X.\]
\end{enumerate}
\endgroup
In 1930 Hausdorff [4] proved that if $A$ is a closed subset of a metrizable space $X$ then any compatible metric $d$ on $A$ can be extended to a compatible metric $\tilde d$ on $X$. Since then the problem of extending metrics has been investigated by several authors [2, 8, 6, 5]. Bacon [1] observed that if $X$ is complete-metrizable and $d$ is complete then $\tilde d$ can be taken to be complete. The aim of this note is to construct two ``canonical'' extensors of metrics preserving, e.g.\ the inequalities of metrics, their point-wise convergence and other properties. The methods used extend those of [7].

Given a metrizable space $X$. By $\Metr(X)$ (resp.\ $\Metr_b(X)$, $\Metr_{tb}(X)$, $\Metr_c(X)$, $\Pm(X)$) we denote the set of all compatible metrics (resp.\ bounded compatible metrics, totally bounded compatible metrics, complete compatible metrics, continuous pseudometrics) on $X$.

If $d_1,d_2\in\Pm(X)$ then
\[\norm{d_1-d_2}=\sup\{|d_1(x,y)-d_2(x,y)|:x,y\in X\}\]
defines a metric (which may take the value $+\infty$) on $\Pm(X)$, moreover $\norm d$ is finite for $d\in\Metr_b(X)$ and $\norm{d_1-d_2}$ may be finite for non-bounded pseudometrics $d_1,d_2\in\Pm(X)$.

A map $\Phi:\Pm(X)\to\Pm(Y)$ is Lipschitz iff there exists $C>0$ such that
\[\norm{\Phi(d_1)-\Phi(d_2)}\leq C\norm{d_1-d_2}\quad\text{for }d_1,d_2\in\Pm(X).\]

\newpage % Original page 286
Let $X$ be a metrizable space and $A$ be a closed subset of $X$. In this note we prove the following two theorems

\par\medskip\noindent\textsc{Theorem 1.} \textit{There exists an extensor $\Phi:\Pm(A)\to\Pm(X)$ satisfying the conditions}
\begin{conditions}{1}
\item $\Phi(d)\restr A=d$ for $d\in\Pm(A)$.
\item $\Phi(\Metr(A))\subset\Metr(X)$.
\item \textit{$\Phi$ is a Lipschitz map.}
\item $\Phi(\Metr_b(A))\subset\Metr_b(X)$.
\item \textit{If $X$ is separable then $\Phi(\Metr_{tb}(A))=\Metr_{tb}(X)$.}
\item \textit{If $X$ is complete-metrizable then $\Phi(\Metr_c(A))-\Metr_c(X)$.}
\end{conditions}

\par\medskip\noindent\textsc{Theorem 2.} \textit{There exists an extensor $\Phi:\Pm(A)\to\Pm(X)$ satisfying the conditions}
\begin{conditions}{2}
\item $\Phi(d)\restr A=d$ for $d\in\Pm(A)$.
\item $\Phi(\Metr(A))\subset\Metr(X)$.
\item \textit{If $d_1\leq d_2$ then $\Phi(d_1)\leq\Phi(d_2)$ for $d_1,d_2\in\Pm(A)$.}
\item $\Phi(\Metr_b(A))\subset\Metr_b(X)$.
\item \textit{If $X$ is separable then $\Phi(\Metr_{tb}(A))\subset\Metr_{tb}(X)$.}
\item \textit{If $X$ is complete-metrizable then $\Phi(\Metr_c(A))\subset\Metr_c(X)$}
\item \textit{If $d_n(x,y)\to d_0(x,y)$ for $x,y\in A$ then}
\[\Phi(d_n)(x,y)\to\Phi(d_0)(x,y)\quad\text{for }x,y\in X.\]
\end{conditions}

\par\medskip\noindent\textbf{1. Proof of Theorem 1.} Consider the locally convex space $E=m(\Pm(A),l^\infty(A))$ of all maps from $\Pm(A)$ into $l^\infty(A)$ equipped with the topology of point-wise convergence. Fix a base point $x_0\in A$ and define a map $g:A\to E$ by
\begin{equation}\tag{1-1}
g_x(d)(a)=d(x,a)-d(x_0,a)\quad\text{for }d\in\Pm(A)\text{ and }a\in A.
\end{equation}
Obviously $g$ is continuous. By Dugundji's theorem [3] there exists a continuous map $\tilde g:X\to E$ such that $\tilde g\restr A=g$ and $\tilde g(X)\subset\conv g(A)$ (the convex hull of $g(A)$ in $E$).

Fix a compatible metric $\rho$ on $X$ which is complete if $X$ is complete-metrizable and totally bounded if $X$ is separable and define $\Phi_1,\Phi_2,\Phi:\Pm(A)\to\Pm(X)$ by
\begin{equation}\tag{1-2}
\Phi_1(d)(x,y)=\norm{g_x(d)-g_y(d)}
\end{equation}
\begin{equation}\tag{1-3}
\Phi_2(d)(x,y)=\sup\{|\theta(d)(x,z)\,\theta(d)(x,A)-\theta(d)(y,z)\,\theta(d)(y,A)|:z\in X\},
\end{equation}
where
\begin{equation}\tag{1-4}
\theta(d)=\min\{1,\Phi_1(d)+\rho\}
\end{equation}
\begin{equation}\tag{1-5}
\Phi(d)=\Phi_1(d)+\Phi_2(d).
\end{equation}
Obviously $\Phi$ satisfies (1-i).

\proofpart{1-ii}Since $\theta(\Pm(A))\subset\Metr(X)$, it suffices to show that $\Phi(d)$ is equivalent to $\theta(d)$ for every $d\in\Metr(A)$.

Suppose that $\theta(d)(x_n,x)\to0$. By (1-4) and (1-3) we have $\Phi_i(d)(x_n,x)\to0$ for $i=1,2$. Thus $\Phi(d)(x_n,x)\to0$ by (1-5).

\newpage % Original page 287
Conversely, assume to the contrary that for some sequence $(x_n)\subset X$ we have $\Phi(d)(x_n,x)\to0$ and
\begin{equation}\tag{1-6}
\theta(d)(x_n,x)\geq\delta>0\quad\text{for every }n\in N.
\end{equation}
Since by (1-3)
\[
\theta(d)(x_n,x)\max\{\theta(d)(x,A),\theta(d)(x_n,A)\}\leq\Phi_2(d)(x_n,x)\leq\Phi(d)(x_n,x)\to0
\]
we infer that $x\in A$ and $\theta(d)(x_n,A)\to0$.

Select a sequence of points $(x'_n)\subset A$ such that
\begin{equation}\tag{1-7}
\theta(d)(x_n,x'_n)\to0
\end{equation}
then $\Phi_1(d)(x_n,x'_n)\to0$, by (1-4).

Since $\Phi(d)(x_n,x)\to0$, from (1-5) we get $\Phi_1(d)(x_n,x)\to0$ and
\begin{equation}\tag{1-8}
d(x'_n,x)=\Phi_1(d)(x'_n,x)\to0.
\end{equation}
Note that $x\in A$, $(x_n)\subset A$ and $d,\theta(d)\restr A\in\Metr(A)$, from (1-4) and (1-8) it follows that $\theta(d)(x'_n,x)\to0$. Combining this with (1-7) we obtain $\theta(d)(x_n,x)\to0$, which contradicts (1-6).

\proofpart{1-iii}From (1-1) we get
\begin{equation}\tag{1-9}
\norm{g_x(d_1)-g_x(d_2)}\leq2\norm{d_1-d_2}\quad\text{for every }x\in A.
\end{equation}
Since $\tilde g(X)\subset\conv g(A)$, for every $x,y\in X$ we have
\[
\begin{aligned}
g_x&=\sum_{i=1}^n\lambda_i g_{x_i},\quad x_i\in A,\quad\lambda_i\geq0\text{ for }i=1,\ldots,n\text{ and }\sum_{i=1}^n\lambda_i=1,\\[5pt]
g_y&=\sum_{j=1}^m\mu_j g_{y_j},\quad y_j\in A,\quad\mu_j\geq0\text{ for }j=1,\ldots,m\text{ and }\sum_{j=1}^m\mu_j=1.
\end{aligned}
\]
Thus from (1-9) and (1-2) we obtain
\[\norm{\Phi_1(d_1)-\Phi_1(d_2)}\leq4\norm{d_1-d_2}.\]
On the other hand, from (1-3) and (1-4) we get
\[
\norm{\Phi_2(d_1)-\Phi_2(d_2)}\leq4\norm{\theta(d_1)-\theta(d_2)}\leq4\norm{\Phi_1(d_1)-\Phi_1(d_2)}\leq16\norm{d_1-d_2}.
\]
Thus by (1-5)
\[\norm{\Phi(d_1)-\Phi(d_2)}\leq20\norm{d_1-d_2}\quad\text{for }d_1,d_2\in\Pm(A)\]
(1-iv) follows from the relation $\tilde g(X)\subset\conv g(A)$.

\proofpart{1-v}Assume that $\rho$ and $d$ are totally bounded. Since
\[\norm{g_x(d)-g_y(d)}=d(x,y)\]
the set
\[g(A)d=\{g(d):a\in A\}\subset l^\infty(A)\]

\newpage % Original page 288
\noindent is totally bounded. Thus $\tilde g(X)d\subset\conv g(A)d$ is also totally bounded. Hence by (1-2) $\Phi_1(d)$ is totally bounded. Consequently by (1-4) $\theta(d)$ is totally bounded. Since $\Phi_2(d)\leq2\theta(d)$ we infer that
\[\Phi(d)=\Phi_1(d)+\Phi_2(d)\]
is totally bounded.

\proofpart{1-vi}By (1-4) if $\rho$ is complete then $\theta(d)$ is also a complete metric for every $d\in\Pm(A)$. Thus from the equivalence of $\Phi(d)$ and $\theta(d)$ it suffices to show that if $\Phi(d)(x_n,x_m)\to0$ then $\theta(d)(x_n,x_m)\to0$ for every $d\in\Metr_c(A)$.

If it is not the case then there exist two subsequences $(a_n)\subset(x_n)$ and $(b_n)\subset(x_n)$ such that
\begin{equation}\tag{1-10}
\theta(d)(a_n,b_n)\geq\delta>0\quad\text{for every }n\in N.
\end{equation}
Since by (1-3)
\[\theta(d)(a_n,b_n)\,\theta(d)(a_n,A)\leq\Phi_2(d)(a_n,b_n)\leq\Phi(d)(a_n,b_n)\to0\]
we infer that there exists a sequence $(y_n)$ of points in $A$ such that
\begin{equation}\tag{1-11}
\theta(d)(a_n,y_n)\to0.
\end{equation}
Thus from (1-4) $\Phi_1(d)(a_n,y_n)\to0$. Since $\Phi_1(d)\restr A=d$ is complete and since
\[
\begin{aligned}
\Phi_1(d)(y_n,y_m)&\leq\Phi_1(d)(y_n,a_n)+\Phi_1(d)(a_n,a_m)+\Phi_1(d)(a_m,y_m)\\
&\leq\Phi_1(d)(y_n,a_n)+\Phi(d)(a_n,a_m)+\Phi_1(d)(a_m,y_m)\to0
\end{aligned}
\]
we infer that there exists a point $x_0\in A$ such that
\[\Phi(d)(y_n,x_0)=\Phi_1(d)(y_n,x_0)=d(y_n,x_0)\to0.\]
Thus $\theta(d)(y_n,x_0)\to0$. Hence from (1-11) we get $\theta(d)(a_n,x_0)\to0$. Consequently $\Phi(d)(a_n,x_0)\to0$. Since $(x_n)$ is a $\Phi(d)$-Cauchy sequence and $(a_n)\subset(x_n)$ we infer that $\Phi(d)(x_n,x_0)\to0$. Thus $\theta(d)(x_n,x_0)\to0$. This contradicts (1-10) and concludes the proof of Theorem 1.

\par\medskip\noindent\textbf{2. Proof of Theorem 2.} For every $d\in\Pm(A)$ let $\Lip(A,d)$ denote the set of all Lipschitz maps $f:(A,d)\to R^1$ vanishing at $a_0\in A$ equipped with the norm
\[\norm f_d=\inf\{C:f\text{ is a }C\text{-Lipschitz map}\}.\]
Let $L(A)$ denote the liner space spanned formally by $A$. For every $d\in\Pm(A)$ and $x=\sum_{i=1}^n\lambda_i a_i\in L(A)$, put
\[
\norm x_d=\sup\left\{\left|\sum_{i=1}^n\lambda_i f(a_i)\right|:f\in\Lip(A,d),\ \norm f_d\leq1\right\}.
\]
Then $L(A)$ becomes a locally convex space. By Dugundji's theorem [3] there exists a continuous map $\varphi:X\to L(A)$ such that $\varphi(a)=a$ for every $a\in A$ and $\varphi(X)\subset\conv(A)\subset L(A)$.

\newpage % Original page 289
Fix a compatible metric $\rho$ on $X$ which is complete if $X$ is complete-metrizable and totally bounded if $X$ is separable and define $\Phi_1,\Phi_2,\Phi:\Pm(A)\to\Pm(X)$ by
\begin{equation}\tag{2-1}
\Phi_1(d)(x,y)=\norm{\varphi(x)-\varphi(y)}_d
\end{equation}
\begin{equation}\tag{2-2}
\Phi_2(d)(x,y)=\sup\{|f(x)-f(y)|:f\in\Lip(X,e),\ f\restr A=0\text{ and }\norm f_e\leq1\},
\end{equation}
where $e=\min\{1,\rho+\Phi_1(d)\}$
\begin{equation}\tag{2-3}
\Phi(d)=\Phi_1(d)+\Phi_2(d).
\end{equation}
It is easy to see that $\Phi$ satisfies conditions (2-i) and (2-iii). The proofs of (2-ii), (2-iv), (2-v) and (2-vi) are similar to the proofs of (1-ii), (1-iv), (1-v) and (1-vi), respectively. Let us check (2-vii).

Assume that $d_n(x,y)\to d_0(x,y)$ for $x,y\in A$. We have to show that $\Phi_i(d_n)(x,y)\to\Phi_i(d_0)(x,y)$ for $x,y\in X$ and $i=1,2$.

Take $x_0,y_0\in X$. Then we have
\[\varphi(x_0)-\varphi(y_0)=\sum_{i=1}^k\lambda_i a_i,\]
where $a_1,a_2,\ldots,a_k\in A$.

\textit{The case $i=1$.} For every $\varepsilon>0$, take $p>\Phi_1(d_0)(x_0,y_0)$ and let $n_0\in N$ be such that
\begin{equation}\tag{2-4}
\left(1-\frac\varepsilon p\right)d_0(a,b)\leq d_n(a,b)\leq\left(1+\frac\varepsilon p\right)d_0(a,b)
\end{equation}
for every $n\geq n_0$ and $a,b\in A_0=\{a_1,\ldots,a_k\}$.

From (2-4) we get for $n\geq n_0$
\[
\begin{aligned}
\Phi_1(d_n)(x_0,y_0)
&=\sup\left\{\left|\sum_{i=1}^k\lambda_i f(a_i)\right|:f\in\Lip(A,d_n):\norm f_{d_n}\leq1\right\}\\[5pt]
&=\sup\left\{\left|\sum_{i=1}^k\lambda_i f(a_i)\right|:f\in\Lip(A_0,d_n),\ \norm f_{d_n}\leq1\right\}\\[5pt]
&\leq\sup\left\{\left|\sum_{i=1}^k\lambda_i f(a_i)\right|:f\in\Lip(A_0,d_0),\ \norm f_{d_0}\leq1+\frac\varepsilon p\right\}\\[5pt]
&=\left(1+\frac\varepsilon p\right)\sup\left\{\left|\sum_{i=1}^k\lambda_i f(a_i)\right|:f\in\Lip(A_0,d_0),\ \norm f_{d_0}\leq1\right\}\\[5pt]
&=\left(1+\frac\varepsilon p\right)\sup\left\{\left|\sum_{i=1}^k\lambda_i f(a_i)\right|:f\in\Lip(A,d_0),\ \norm f_{d_0}\leq1\right\}\\[5pt]
&=\left(1+\frac\varepsilon p\right)\Phi_1(d_0)(x_0,y_0)<\Phi_1(d_0)(x_0,y_0)+\varepsilon.
\end{aligned}
\]
Similarly,
\[\Phi_1(d_n)(x_0,y_0)>\Phi_1(d_0)(x_0,y_0)-\varepsilon\quad\text{for every }n\geq n_0.\]

\newpage % Original page 290
\textit{The case $i=2$.} Let $e_n(x,y)=\rho(x,y)+\Phi_1(d_n)(x,y)$. By the case $i=1$ we get
\[e_n(x,y)\to e_0(x,y)\quad\text{for every }x,y\in X.\]
For every $\varepsilon>0$, take an $n_0\in N$ such that
\begin{equation}\tag{2-5}
\left(1-\frac\varepsilon p\right)e_0(x,y)\leq e_n(x,y)\leq\left(1+\frac\varepsilon p\right)e_0(x_0,y_0)\quad\text{for every }n\geq n_0,
\end{equation}
where $p>\Phi_2(d_0)(x_0,y_0)$.

Thus for every $n\geq n_0$ we have
\[
\begin{aligned}
\Phi_2(d_n)(x_0,y_0)
&=\sup\{|f(x_0)-f(y_0)|:f\in\Lip(X,e_n),\ f\restr A=0\text{ and }\norm f_{e_n}\leq1\}\\[6pt]
&\leq\sup\left\{|f(x_0)-f(y_0)|:f\in\Lip(X,e_0),\ f\restr A=0\text{ and }\norm f_{e_0}\leq1+\frac\varepsilon p\right\}\\[6pt]
&=\left(1+\frac\varepsilon p\right)\sup\{|f(x_0)-f(y_0)|:f\in\Lip(X,e_0),\ f\restr A=0\text{ and }\norm f_{e_0}\leq1\}\\[6pt]
&=\left(1+\frac\varepsilon p\right)\Phi_2(d_0)(x_0,y_0)<\Phi_2(d_0)(x_0,y_0)+\varepsilon.
\end{aligned}
\]
Similarly,
\[\Phi_2(d_n)(x_0,y_0)>\Phi_2(d_0)(x_0,y_0)-\varepsilon\quad\text{for every }n\geq n_0.\]
This completes the proof of Theorem 2.

\medskip
The authors wish to express their deep gratitude to H. Toru\'nczyk for his kind interest during the preparation of this note.

Their thanks are due also to J. Luukkainen for many valuable remarks.

\medskip
{\footnotesize\noindent INSTITUTE OF MATHEMATICS, POLISH ACADEMY OF SCIENCES, \'SNIADECKICH 8, 00-656 WARSAW (INSTYTUT MATEMATYCZNY, PAN).\par}

\medskip\noindent REFERENCES
\begin{enumerate}[label={[\arabic*]},leftmargin=2em,itemsep=3pt,topsep=7pt,parsep=0pt]
\item P. Bacon, \textit{Extending a complete metric}, Amer. Math. Monthly, \textbf{75} (1968), 642--643.
\item R. H. Bing, \textit{Extending a metric}, Duke Math. J., \textbf{14} (1947), 511--519.
\item J. Dugundji, \textit{An extension of Tietze's theorem}, Pacific J. Math., \textbf{1} (1951), 353--367.
\item F. Hausdorff, \textit{Erweiterung einer Homomorphie}, Fund. Math., \textbf{16} (1930), 353--360.
\item J. Luukkainen, \textit{Extension of spaces, maps and metrics in Lipschitz topology}, Ann. Acad. Sci. Fenn. Ser. A1 Math. Dissertations, \textbf{17} (1978), 1--62.
\item Nguyen To Nhu, \textit{Extending metrics uniformly}, Coll. Math., \textbf{42} (1) [to appear].
\item ---\,, \textit{$P^\lambda$-spaces and $L^\lambda$-spaces}, ibid., \textbf{41} (1979), 67--71.
\item H. Toru\'nczyk, \textit{A short proof of Hausdorff's theorem on extending metrics}, Fund. Math., \textbf{77} (1972), 191--193.
\end{enumerate}

\newpage % Original page 291
\begin{otherlanguage}{russian}
\noindent\textbf{Нгуен Ван \foreignlanguage{english}{Khue}, Нгуен То Ну, О двух экстенсорах метрик.}

Через $\Metr(Z)$ (соответственно, $\Metr_c(Z)$) мы обозначаем семейство всех компактизуемы метрик (соответственно, вполне компактизуемых метрик) на данном метризуемом пространстве $Z$. Пусть $A$ замкнутое подмножество метризуемого пространства $X$. В данной работе строятся два экстенсора метрик $\Phi_1,\Phi:\Metr(A)\to\Metr(X)$, удовлетворяющих условиям:
\begin{enumerate}[label=(\roman*),leftmargin=3.5em,itemsep=3pt,topsep=4pt]
\item $\Phi_i(d)\restr A=d$, где $d\in\Metr(A)$, $i=1,2$;
\item если $X$ вполне метризуемо, то $\Phi_i(\Metr_c(A))\subset\Metr_c(X)$ $i=1,2$;
\item $\Phi_1$ есть Липшицово отображение;
\item если $d_1\leq d$, то $\Phi(d_1)\leq\Phi(d)$ где $d_1,d\in\Metr(A)$;
\item если $d_n(x,y)$, где $x,y\in A$, то $\Phi(d_n)(x,y)\quad\Phi(d)(x,y)$, где $x,y\in X$.
\end{enumerate}
\end{otherlanguage}
\end{document}
